\chapter{مدارات الشظايا والجزيئات المتعددة الذرات}\label{ch:b2:fragment-orbitals}

تُظهر بنية لويس للماء رابطتين O--H متكافئتين وزوجين حرين
متكافئين. ومع ذلك فحين يُؤيَّن الماء بالأشعة فوق البنفسجية تخرج
الإلكترونات بأربع طاقات مختلفة، لا باثنتين؛ والميثان، بروابطه C--H الأربع
المتكافئة، يُظهر طاقتين. فالإلكترونات في الجزيء لا تجلس في روابط
بين أزواج من الذرات: بل تشغل مدارات ممتدة على الجزيء كله.
يبني هذا الفصل هذه المدارات من \omterm{def:b2:fragment-orbitals:fragment}{مدارات الشظايا}، ويفسر
لماذا يكون الماء منحنيًا والإيثين مستويًا، ولماذا يفضّل الكاتيون الكربوني أن تحيط به
مجموعات ألكيل.

\begin{recall}
\Cref{ch:b2:diatomic-mos}: طريقة \omterm{def:b2:diatomic-mos:lcao}{LCAO}، والتداخل والتناظر، ومستويان متآثران،
\omterm{def:b2:diatomic-mos:bonding}{والمدارات الرابطة} والمضادة للربط. ومجلد السنة الأولى: أشكال VSEPR، والتهجين
وصفًا، والرنين، والكاتيونات الكربونية وترتيب استقرارها (وقد أُعلن هناك
تآثر روابط C--H مع المدار الفارغ).
\end{recall}

\section{طريقة الشظايا}

\begin{definition}[الشظايا]\label{def:b2:fragment-orbitals:fragment}
\emph{الشظية}\index{شظية} مجموعة من ذرات الجزيء تُعتبر
وحدها (ذرة، أو زوج من الهيدروجينات، أو مجموعة \ce{CH2})، و\emph{مدار
الشظية}\index{مدار الشظية} أحد مداراتها. وتُبنى مدارات
الجزيء بجعل مدارات شظيتين تتآثر، اثنين
اثنين، كما تتآثر المدارات الذرية في جزيء ثنائي الذرة.
\end{definition}

\begin{definition}[التركيب المتكيف مع التناظر]\label{def:b2:fragment-orbitals:symmetry-adapted}
\emph{التركيب المتكيف مع التناظر}\index{تركيب متكيف مع التناظر} لمدارات ذرية
متكافئة تركيب تتركه كل عملية تناظر في
الجزيء (انعكاس في مستوي مرآة، أو دوران حول محور) إما
دون تغيير وإما تحوّله إلى معاكسه، أو، في المجموعات المتساوية الطاقة، تحوّله
إلى تراكيب من المجموعة نفسها.
\end{definition}

\begin{proposition}[هيدروجينان]\label{prop:b2:fragment-orbitals:h2-pair}
من أجل هيدروجينين متكافئين مداراهما 1s هما $h_1$ و$h_2$، يكون
\omterm{def:b2:fragment-orbitals:symmetry-adapted}{التركيبان المتكيفان مع التناظر} $h_1 + h_2$ و$h_1 - h_2$: فالأول لا
يتغير، والثاني يغيّر إشارته، تحت أي عملية تبادل بين
الذرتين.
\end{proposition}

\begin{proof}
العملية التي تبادل بين الذرتين تحوّل $h_1$ إلى $h_2$ و$h_2$ إلى
$h_1$: فيبقى $h_1 + h_2$ دون تغيير ويصير $h_1 - h_2$ مساويًا $h_2 - h_1 = -(h_1 - h_2)$.
\end{proof}

\begin{proposition}[ثلاثة هيدروجينات وأربعة]\label{prop:b2:fragment-orbitals:h4}
من أجل أربعة هيدروجينات عند رؤوس رباعي أوجه، تكون التراكيب تركيبًا
تام التناظر $h_1 + h_2 + h_3 + h_4$ ومجموعةً من ثلاثة تراكيب متساوية
الطاقة لكل منها مستوٍ عقدي واحد، على شكل مدارات p. ومن أجل ثلاثة
هيدروجينات على مثلث (كما في الأمونيا)، تركيب متناظر واحد وزوج
متساوي الطاقة.
\end{proposition}

\admitted

تأتي أشكال التراكيب المتساوية الطاقة، وأسماؤها (a، e، t)، من
نظرية الزمر، التي تُطوَّر في مجلد السنة الثالثة؛ وهنا ليست الأسماء إلا وسومًا.

\begin{proposition}[لا يتآثر إلا المتماثل مع المتماثل]\label{prop:b2:fragment-orbitals:interaction-rules}
لا يتآثر مدار \omterm{def:b2:fragment-orbitals:fragment}{شظية} إلا مع \omterm{def:b2:fragment-orbitals:fragment}{مدارات الشظية} الأخرى التي
تسلك السلوك نفسه تحت كل عملية تناظر في الجزيء؛ وتزداد
شدة التآثر مع التداخل وتنقص مع الفجوة
الطاقية.
\end{proposition}

\begin{proof}
إذا سلك مداران سلوكين مختلفين تحت انعكاس ما، كان أحدهما متناظرًا 
والآخر ضد متناظر بالنسبة إلى ذلك المستوي، فينعدم تداخلهما وتكامل
رنينهما (\cref{prop:b2:diatomic-mos:symmetry}). والتعلق بالتداخل
وبالفجوة هو ما في \cref{thm:b2:diatomic-mos:heteronuclear}.
\end{proof}

\begin{method}[مخطط الشظايا]\label{met:b2:fragment-orbitals:fragment-diagram}
\begin{enumerate}
  \item اقسم الجزيء إلى شظيتين، عادة ذرة مركزية ومجموعة
        الذرات المحيطة بها.
  \item ابنِ \omterm{def:b2:fragment-orbitals:symmetry-adapted}{التراكيب المتكيفة مع التناظر} للمدارات الخارجية.
  \item صنّف \omterm{def:b2:fragment-orbitals:fragment}{مدارات الشظيتين} وفق سلوكها تحت مستويات المرآة
        ومحاور الجزيء.
  \item اجعل كل مدار يتآثر مع المدارات ذات السلوك نفسه، بقوة أكبر
        حين تتقارب طاقاتها؛ وتبقى المدارات المتبقية غير رابطة.
  \item املأ بإلكترونات التكافؤ واقرأ: أعلى مستوى مشغول،
        والطابع الرابط، وتفضيلات الشكل.
\end{enumerate}
\end{method}

\section{الماء}

ضع الماء في المستوي $yz$، والمحور $z$ منصِّف للزاوية H--O--H. يحوي
مستويا مرآة المحور $z$: المستوي الجزيئي $yz$ والمستوي $xz$
العمودي عليه. وبالنسبة إلى المستوي $xz$، الذي يبادل بين
الهيدروجينين، يكون $h_1 + h_2$ متناظرًا و$h_1 - h_2$ ضد متناظر. وعلى الأكسجين، يكون 2s 
و$2\mathrm p_z$ متناظرين بالنسبة إلى المستويين كليهما؛ و$2\mathrm p_y$ (في
المستوي الجزيئي، عمودي على $z$) ضد متناظر بالنسبة إلى $xz$؛
و$2\mathrm p_x$ ضد متناظر بالنسبة إلى المستوي الجزيئي، وهو ما لا يكونه أي
تركيب هيدروجيني.

وهكذا يمتزج $h_1 + h_2$ مع 2s و$2\mathrm p_z$، و$h_1 - h_2$ مع $2\mathrm p_y$، 
ويبقى $2\mathrm p_x$ غير رابط. فتنتج ستة مدارات، يملأ أربعة منها إلكترونات
التكافؤ الثمانية.

\begin{omfigure}
\begin{tikzpicture}[font=\small]
  % O fragment (left)
  \pic (os) at (0,-2.4) {omlevel={ud}{0.7}}; \node[left] at (-0.45,-2.4) {2s};
  \pic (op1) at (-0.45,0) {omlevel={ud}{0.4}}; \pic (op2) at (0,0) {omlevel={u}{0.4}}; \pic (op3) at (0.45,0) {omlevel={u}{0.4}};
  \node[left] at (-0.7,0) {2p};
  \node at (0,-3.2) {O};
  % H2 fragment (right)
  \pic (hp) at (6,0.9) {omlevel={u}{0.7}}; \node[right] at (6.4,0.9) {$h_1 + h_2$};
  \pic (hm) at (6,1.8) {omlevel={u}{0.7}}; \node[right] at (6.4,1.8) {$h_1 - h_2$};
  \node[anchor=north west] at (6.4,0.6) {\foreignlanguage{arabic}{الشظية \ce{H2}}};
  % MOs
  \pic (m1) at (3,-2.9) {omlevel={ud}{0.7}};  \node[omsymlabel, below=0.17cm, inner sep=0.5pt] at (m1-c) {$2a_1$};
  \pic (m2) at (3,-1.55) {omlevel={ud}{0.7}};  \node[omsymlabel, below=0.17cm, inner sep=0.5pt] at (m2-c) {$1b_2$};
  \pic (m3) at (3,-0.78) {omlevel={ud}{0.7}};  \node[omsymlabel, below=0.17cm, inner sep=0.5pt] at (m3-c) {$3a_1$};
  \pic (m4) at (3,0) {omlevel={ud}{0.7}};  \node[omsymlabel, below=0.17cm, inner sep=0.5pt] at (m4-c) {$1b_1$};
  \pic (m5) at (3,2.4) {omlevel={empty}{0.7}}; \node[omsymlabel, below=0.17cm, inner sep=0.5pt] at (m5-c) {$4a_1$};
  \pic (m6) at (3,3.1) {omlevel={empty}{0.7}}; \node[omsymlabel, below=0.17cm, inner sep=0.5pt] at (m6-c) {$2b_2$};
  \draw[omcorr, densely dashed] (os-r) -- (m1-l) (op3-r) -- (m2-l) (op3-r) -- (m3-l) (op3-r) -- (m4-l) (op3-r) -- (m5-l) (op3-r) -- (m6-l);
  \draw[omcorr, densely dashed] (hp-l) -- (m1-r) (hm-l) -- (m2-r) (hp-l) -- (m3-r) (hp-l) -- (m5-r) (hm-l) -- (m6-r);
\end{tikzpicture}

{\small مخطط الشظايا للماء، رسم تخطيطي. يمتزج التركيب المتناظر
للهيدروجينين مع المدارين 2s و$2\mathrm p_z$ للأكسجين (مدارات $a_1$)،
والتركيب ضد المتناظر مع $2\mathrm p_y$ ($b_2$)؛ ويبقى $2\mathrm p_x$، العمودي على
المستوي الجزيئي، غير رابط ($1b_1$، عند مستوى 2p للأكسجين). أربعة مستويات تكافؤ
ممتلئة، وأربع حزم إلكترونات ضوئية؛ وأعلاها، $1b_1$، يتأين عند
\qty{12.62}{eV}. والوسوم
$a_1$ و$b_1$ و$b_2$ أسماء (من نظرية الزمر، في مجلد السنة الثالثة).}
% ledger: ie:H2O
\end{omfigure}

\begin{proposition}[لماذا يكون الماء منحنيًا]\label{prop:b2:fragment-orbitals:bent-water}
في H--O--H خطي، يكون مدارا p للأكسجين العموديان على المحور
غير رابطين. والانحناء يسمح لأحدهما، الواقع في المستوي الجزيئي، بالتداخل
مع $h_1 + h_2$ وبالامتزاج مع المدار 2s: فينخفض ذلك المستوى الممتلئ، بأكثر مما
ترتفع المستويات الممتلئة الأخرى. ومع ثمانية إلكترونات تكافؤ يكون الجزيء المنحني
هو الأكثر استقرارًا.
\end{proposition}

\begin{proof}[تعليل]
في الجزيء الخطي ($z$ على طول H--O--H)، يكون تداخل $h_1 + h_2$ معدومًا مع
$2\mathrm p_x$ و$2\mathrm p_y$ (فكل منهما ضد متناظر بالنسبة إلى مستوٍ
يحوي المحور، والتركيب متناظر)؛ فهما زوجان غير رابطين متساويا
الطاقة. وعند الانحناء في المستوي $yz$، يتجه $2\mathrm p_y$ جزئيًا نحو
الهيدروجينين: فيزداد تداخله مع $h_1 + h_2$ انطلاقًا من الصفر، ويستقر المستوى الممتلئ
الذي يحمله (المستوى $3a_1$ للجزيء المنحني). ويفقد المستوى الرابط
$1b_2$ قليلًا من التداخل فيرتفع قليلًا. والحصيلة انخفاض،
أكبر ما يكون حين يكون المستوى الهابط ممتلئًا، كما هو الحال مع ثمانية إلكترونات.
أما الجزيء الذي ليس له إلا أربعة إلكترونات تكافؤ، مثل \ce{BeH2}، فيكون ذلك المستوى فيه
فارغًا ويكون خطيًا.
\end{proof}

\section{الميثان والأمونيا وأطياف الإلكترونات الضوئية}

في الميثان تعطي مدارات الهيدروجين الأربعة التركيب المتناظر، الذي
يوافق المدار 2s للكربون، ومجموعة متساوية الطاقة من ثلاثة، توافق
مداراته 2p الثلاثة. فتُملأ أربعة \omterm{def:b2:diatomic-mos:bonding}{مدارات رابطة}: واحد منخفض ($a_1$) وثلاثة
متساوية الطاقة عند طاقة أعلى ($t_2$). أما الأمونيا، بهيدروجيناتها الثلاثة، فتحتفظ
بمدار ممتلئ على النيتروجين يتجه بعيدًا عن الهيدروجينات، هو زوجها الحر،
أعلى مستوى مشغول.

\begin{omfigure}
\begin{tikzpicture}[font=\small]
  \pic (cs) at (0,-1.6) {omlevel={ud}{0.7}}; \node[left] at (-0.45,-1.6) {2s};
  \pic (cp1) at (-0.45,0) {omlevel={u}{0.4}}; \pic (cp2) at (0,0) {omlevel={u}{0.4}}; \pic (cp3) at (0.45,0) {omlevel={empty}{0.4}};
  \node[left] at (-0.7,0) {2p};
  \node at (0,-2.5) {C};
  \pic (ha) at (6,-0.8) {omlevel={ud}{0.7}}; \node[right] at (6.4,-0.8) {$a_1$};
  \pic (ht1) at (5.55,0.4) {omlevel={u}{0.4}}; \pic (ht2) at (6,0.4) {omlevel={u}{0.4}}; \pic (ht3) at (6.45,0.4) {omlevel={empty}{0.4}};
  \node[right] at (6.7,0.4) {$t_2$};
  \node at (6,-1.6) {\foreignlanguage{arabic}{الشظية \ce{H4}}};
  \pic (m1) at (3,-2.4) {omlevel={ud}{0.7}}; \node[omsymlabel, below=0.17cm, inner sep=0.5pt] at (m1-c) {$1a_1$};
  \pic (m2a) at (2.55,-0.9) {omlevel={ud}{0.4}}; \pic (m2b) at (3,-0.9) {omlevel={ud}{0.4}}; \pic (m2c) at (3.45,-0.9) {omlevel={ud}{0.4}};
  \node[omsymlabel, below=0.17cm, inner sep=0.5pt] at (m2b-c) {$1t_2$};
  \pic (m3) at (3,1.8) {omlevel={empty}{0.7}}; \node[omsymlabel, below=0.17cm, inner sep=0.5pt] at (m3-c) {$2a_1$};
  \pic (m4a) at (2.55,2.5) {omlevel={empty}{0.4}}; \pic (m4b) at (3,2.5) {omlevel={empty}{0.4}}; \pic (m4c) at (3.45,2.5) {omlevel={empty}{0.4}};
  \node[omsymlabel, below=0.17cm, inner sep=0.5pt] at (m4b-c) {$2t_2$};
  \draw[omcorr, densely dashed] (cs-r) -- (m1-l) (cp3-r) -- (m2a-l) (cs-r) -- (m3-l) (cp3-r) -- (m4a-l);
  \draw[omcorr, densely dashed] (ha-l) -- (m1-r) (ht1-l) -- (m2c-r) (ha-l) -- (m3-r) (ht1-l) -- (m4c-r);
\end{tikzpicture}

{\small مخطط الشظايا للميثان، رسم تخطيطي. يوافق التركيب المتناظر
للهيدروجينات الأربعة المدار 2s للكربون، وتوافق التراكيب الثلاثة المتساوية الطاقة مداراته 2p:
مستويان ممتلئان، ومن ثم حزمتا إلكترونات ضوئية، العليا منهما
(التي تتأين أولًا، عند \qty{12.61}{eV}) تحمل ستة إلكترونات.}
% ledger: ie:CH4
\end{omfigure}

\begin{definition}[مطيافية الإلكترونات الضوئية]\label{def:b2:fragment-orbitals:photoelectron}
\emph{مطيافية الإلكترونات الضوئية}\index{مطيافية الإلكترونات الضوئية} تقيس
الطاقات الحركية للإلكترونات المقذوفة من الجزيئات بإشعاع أحادي اللون
معلوم الطاقة؛ وكل حزمة من الطيف توافق نزع
إلكترون من مستوى مشغول، ويعطي موقعها طاقة تأين
ذلك المستوى.
\end{definition}

\begin{proposition}[تقريب كوبمانز]\label{prop:b2:fragment-orbitals:koopmans}
الطاقة اللازمة لنزع إلكترون من مدار مشغول تساوي
تقريبًا سالب طاقة ذلك المدار، $I \approx -\varepsilon$.
\end{proposition}

\admitted

يهمل التقريب استرخاء الإلكترونات الأخرى بعد
التأين؛ ويُعلَّل في مجلد السنة الثالثة. وهو يحوّل طيف الإلكترونات الضوئية
إلى مخطط مدارات: فالماء يُظهر أربع حزم تكافؤ، والميثان اثنتين،
بما يتفق مع مخططات الشظايا --- وخلافًا لصورة زوجين حرين متكافئين
ورابطتين متكافئتين في الماء، التي تصف الكثافة الإلكترونية الكلية نفسها
بطريقة أخرى لكنها لا تعطي طاقات الإلكترونات
المنزوعة.

\section{الإيثين}

يمكن بناء الإيثين من شظيتين \ce{CH2}. لكل \ce{CH2}، إلى جانب
مداري رابطتيه C--H، مدار يتجه بعيدًا عن الهيدروجينين في
المستوي (يكوّن رابطة $\sigma$ بين C وC مع شريكه) ومدار p
عمودي على المستوي. ويتركب مدارا p في مدار $\pi$ ممتلئ ومدار
$\pi^*$ فارغ.

\begin{omfigure}
\begin{tikzpicture}[font=\small]
  \begin{scope}
    \pic at (0,0) {omp={0.85}{90}}; \pic at (1.2,0) {omp={0.85}{90}};
    \draw[thick] (0,0) -- (1.2,0);
    \draw[thick] (0,0) -- (-0.5,-0.45) (0,0) -- (-0.55,0.35) (1.2,0) -- (1.7,-0.45) (1.2,0) -- (1.75,0.35);
    \node at (0.6,-1.1) {\shortstack{\foreignlanguage{arabic}{$\pi$ (ممتلئ)}}};
  \end{scope}
  \begin{scope}[xshift=3.6cm]
    \pic at (0,0) {omp={0.85}{90}}; \pic at (1.2,0) {omp={-0.85}{90}};
    \draw[thick] (0,0) -- (1.2,0);
    \draw[thick] (0,0) -- (-0.5,-0.45) (0,0) -- (-0.55,0.35) (1.2,0) -- (1.7,-0.45) (1.2,0) -- (1.75,0.35);
    \node at (0.6,-1.1) {\shortstack{\foreignlanguage{arabic}{$\pi^*$ (فارغ)}}};
  \end{scope}
  \begin{scope}[xshift=7.4cm]
    \pic at (0,0) {omp={0.85}{90}}; \pic at (1.2,0) {omp={0.85}{0}};
    \draw[thick] (0,0) -- (1.2,0);
    \node at (0.6,-1.1) {\foreignlanguage{arabic}{ملتوٍ بزاوية $90^\circ$}};
    \node[font=\scriptsize, align=center] at (0.6,1.3) {تداخل معدوم};
  \end{scope}
\end{tikzpicture}

{\small نظام $\pi$ في الإيثين، رسم تخطيطي (الجزيء منظور من حافته، 
والهيدروجينات في المستوي). يتركب مدارا p العموديان على المستوي
بتوافق الطور ($\pi$، ممتلئ) وبتعاكسه ($\pi^*$، فارغ). ولَيّ إحدى مجموعتي \ce{CH2}
بزاوية $90^\circ$ يجعل مداري p متعامدين: فينعدم تداخلهما، ومعه
رابطة $\pi$.}
\end{omfigure}

تُبقي رابطة $\pi$ الإيثين مستويًا: فلَيّ أحد طرفيه يُبعد مداري p أحدهما عن
الآخر، وينخفض تداخلهما كجيب تمام زاوية اللي، وعند $90^\circ$
تزول رابطة $\pi$. ويجب توفير طاقة رابطة $\pi$ لليّ
الجزيء، ولهذا لا يتحول المتماكبان cis وtrans للألكينات أحدهما إلى الآخر في
درجة حرارة الغرفة، بينما يكون الدوران حول رابطة أحادية حرًا.

\section{فرط الاقتران}

\begin{definition}[فرط الاقتران]\label{def:b2:fragment-orbitals:hyperconjugation}
\emph{فرط الاقتران}\index{فرط الاقتران} هو تآثر مدار $\sigma$
ممتلئ لرابطة C--H أو C--C مع مدار p أو $\pi^*$ مجاور فارغ أو ممتلئ
جزئيًا موازٍ له.
\end{definition}

\begin{proposition}[مجموعات الألكيل تثبّت الكاتيونات الكربونية]\label{prop:b2:fragment-orbitals:hyperconjugation}
يستقر الكاتيون الكربوني بكل رابطة C--H أو C--C على كربون مجاور
تستطيع أن تصطف مع مداره p الفارغ: فكلما كثرت هذه الروابط، ازداد استقرار
الكاتيون، ومن هنا الترتيب \babelsublr{\foreignlanguage{arabic}{ثالثي} > \foreignlanguage{arabic}{ثانوي} > \foreignlanguage{arabic}{أولي} > \foreignlanguage{arabic}{ميثيل}}.
\end{proposition}

\begin{proof}[تعليل]
يكوّن المدار p الفارغ للكربون الكاتيوني ومدار $\sigma_{\text{C--H}}$ ممتلئ
مجاور له نظامًا من مستويين بفجوة كبيرة $\Delta$ واقتران صغير
$\beta$ (غير معدوم إلا إذا كانت الرابطة موازية تقريبًا للمدار p). وحسب
\cref{thm:b2:diatomic-mos:heteronuclear}، ينخفض المستوى $\sigma$ الممتلئ بنحو
$\beta^2/\Delta$، ويكسب الإلكترونان ضعف ذلك؛ ويرتفع المستوى الفارغ، دون
كلفة. وكل رابطة مصطفة على نحو مناسب تضيف استقرارها الخاص؛ ولمجموعة الميثيل
دائمًا رابطة C--H شبه مصطفة أيًّا كان دورانها.
\end{proof}

\begin{omfigure}
\begin{tikzpicture}[font=\small]
  \pic at (0,0) {omp={0.9}{90}};
  \node[circle, fill=black, inner sep=1.2pt] at (0,0) {};
  \node[font=\scriptsize, anchor=north west] at (0.08,-0.05) {C$^+$};
  \draw[thick] (0,0) -- (1.4,0);
  \node[circle, fill=black, inner sep=1.2pt] at (1.4,0) {};
  \draw[thick] (1.4,0) -- (1.75,0.85);
  \pic at (1.58,0.42) {omlobe={0.75}{68}};
  \pic at (1.58,0.42) {omlobe={0.45}{248}};
  \node[font=\scriptsize] at (1.95,1.0) {H};
  \draw[thick] (1.4,0) -- (2.0,-0.45) (1.4,0) -- (1.1,-0.6);
  \draw[<->, omProp, thick] (0.35,0.8) .. controls (0.8,1.2) and (1.2,1.2) .. (1.5,0.95);
  \node[font=\scriptsize, align=left, anchor=west] at (2.4,0.4) {\foreignlanguage{arabic}{فص $\sigma_{\text{C--H}}$ ممتلئ موازٍ}\\ \foreignlanguage{arabic}{للمدار \babelsublr{p} الفارغ}};
  % level scheme
  \begin{scope}[xshift=7.2cm, yshift=-0.4cm]
    \pic (p) at (0,1.4) {omlevel={empty}{0.7}}; \node[left] at (-0.4,1.4) {\foreignlanguage{arabic}{\babelsublr{p} فارغ}};
    \pic (s) at (2.2,-0.8) {omlevel={ud}{0.7}}; \node[right] at (2.6,-0.8) {$\sigma_{\text{C--H}}$};
    \pic (lo) at (1.1,-1.2) {omlevel={ud}{0.7}};
    \pic (hi) at (1.1,1.7) {omlevel={empty}{0.7}};
    \draw[omcorr, densely dashed] (p-r) -- (lo-l) (p-r) -- (hi-l) (s-l) -- (lo-r) (s-l) -- (hi-r);
    \node[font=\scriptsize, anchor=north] at (1.1,-1.35) {\foreignlanguage{arabic}{استقرار بمقدار $\approx \beta^2/\Delta$}};
  \end{scope}
\end{tikzpicture}

{\small \omterm{def:b2:fragment-orbitals:hyperconjugation}{فرط الاقتران} في كاتيون كربوني، رسم تخطيطي. \omterm{def:b2:diatomic-mos:bonding}{مدار رابط} C--H ممتلئ
على الكربون المجاور، موازٍ للمدار p الفارغ، يمنحه جزءًا من
كثافته الإلكترونية؛ فيستقر المستوى الممتلئ كما في كل تآثر بين
مستويين مختلفي الطاقة.}
\end{omfigure}

\section{تمارين}

\begin{exercise}[$\star$]\label{exo:b2:fragment-orbitals:1}
اكتب \omterm{def:b2:fragment-orbitals:symmetry-adapted}{التركيبين المتكيفين مع التناظر} لمداري 1s لهيدروجيني
الماء وبيّن أيهما يغيّر إشارته في المستوي الذي ينصّف الزاوية H--O--H
عموديًا على الجزيء.
\end{exercise}

\begin{exercise}[$\star$]\label{exo:b2:fragment-orbitals:2}
كم \omterm{def:b2:diatomic-mos:lcao}{مدارًا جزيئيًا} تكافؤيًا للماء، وكم إلكترون تكافؤ،
وكم مدارًا تكافؤيًا ممتلئًا؟
\end{exercise}

\begin{exercise}[$\star$]\label{exo:b2:fragment-orbitals:3}
لماذا يُظهر طيف الإلكترونات الضوئية للميثان حزمتي تكافؤ، إحداهما
أشد بثلاث مرات من الأخرى؟
\end{exercise}

\begin{exercise}[$\star$]\label{exo:b2:fragment-orbitals:4}
لماذا تقع ذرات الإيثين الست في مستوٍ واحد؟
\end{exercise}

\begin{exercise}[$\star\star$]\label{exo:b2:fragment-orbitals:5}
طاقة التأين الأولى للأمونيا \qty{10.07}{eV}، وللميثان
\qty{12.61}{eV}. من أي مدار يُنزع الإلكترون في كل حالة، ولماذا تكون الأمونيا
أسهل تأينًا؟
\end{exercise}

\begin{exercise}[$\star\star$]\label{exo:b2:fragment-orbitals:6}
قارن المستويات المشغولة للماء الخطي والمنحني، وفسّر لماذا يكون جزيء
مثل \ce{BeH2}، بإلكترونات تكافؤه الأربعة، خطيًا.
\end{exercise}

\begin{exercise}[$\star\star$]\label{exo:b2:fragment-orbitals:7}
فسّر، بنتيجة المستويين، ترتيب استقرار الكاتيونات الكربونية
\ce{CH3+} و\ce{CH3CH2+} و\ce{(CH3)2CH+} و\ce{(CH3)3C+}.
\end{exercise}

\begin{exercise}[$\star\star$]\label{exo:b2:fragment-orbitals:8}
تداخل مداري p ملتويين بزاوية $\theta$ حول الرابطة
يتناسب مع $\cos\theta$. ومع \omterm{def:b2:diatomic-mos:integrals}{تكامل رنين} يتناسب مع
التداخل، كيف يتعلق استقرار $\pi$ بالزاوية $\theta$؟ وعند أي زاوية
ينخفض إلى النصف؟
\end{exercise}

\begin{exercise}[$\star\star$]\label{exo:b2:fragment-orbitals:9}
البوران \ce{BH3} مستوٍ، والأمونيا هرمية. مستعملًا \omterm{def:b2:fragment-orbitals:fragment}{مدارات الشظايا} لذرة
مركزية وثلاثة هيدروجينات، فسّر الفرق بالإلكترونات التي يملكها كل منهما.
\end{exercise}

\begin{exercise}[$\star\star\star$]\label{exo:b2:fragment-orbitals:10}
الأيون \ce{H3+} مثلث متساوي الأضلاع. مع $\alpha$ و$\beta$ لمدارات 1s
وإهمال التداخل، أوجد طاقاته المدارية الثلاث (للتركيب المتناظر
الطاقة $\alpha + 2\beta$) وفسّر لماذا يكفي إلكترونان
لربطه.
\end{exercise}

\begin{exercise}[$\star\star\star$]\label{exo:b2:fragment-orbitals:11}
لنظام الأليل \ce{CH2=CH-CH2} ثلاثة مدارات p متوازية. خمّن أشكال
مداراته $\pi$ الثلاثة (الإشارات على كل كربون) من عدد العقد، وأيها
أعلى مستوى مشغول في أنيون الأليل.
\end{exercise}

\begin{exercise}[$\star\star\star$]\label{exo:b2:fragment-orbitals:12}
ابنِ نظام $\pi$ لثنائي أكسيد الكربون من مدارات 2p العمودية على
المحور O--C--O: كم مدار $\pi$، وأيها رابط، وأيها غير رابط، وأيها
مضاد للربط، وكم إلكترونًا تحمل؟
\end{exercise}

\section{مسألة: شكل الماء}

\begin{problem}[{مسألة نهاية الأسبوع --- الشظيتان الأكسجين و\ce{H2}، ومدارات
H--O--H خطي، والمستويات التي تتحرك حين ينحني، وطيف الإلكترونات
الضوئية مع تقريب كوبمانز}]
\label{pb:b2:fragment-orbitals:1}
المعطيات: هندسة الماء، $r(\ce{O-H}) = \qty{95.8}{pm}$، والزاوية
$104.48^\circ$؛ وطاقات التأين الأولى: \ce{H2O} \qty{12.62}{eV}، \ce{O}
\qty{13.62}{eV}. يقع الماء في المستوي $yz$، و$z$ على طول المنصِّف.
% ledger: geo:H2O.rOH, geo:H2O.aHOH, ie:H2O, ie:O

\medskip
\noindent\textbf{الجزء الأول --- الشظايا.}
\begin{enumerate}
  \item كم إلكترون تكافؤ للماء، ومن أي الذرات؟
  \item اكتب \omterm{def:b2:fragment-orbitals:symmetry-adapted}{التركيبين المتكيفين مع التناظر} لمداري 1s الهيدروجينيين.
  \item أي مدارات الأكسجين متناظرة بالنسبة إلى مستويي المرآة كليهما؟
  \item أي مدار للأكسجين ضد متناظر بالنسبة إلى المستوي $xz$ وحده؟
  \item أي مدار للأكسجين لا شريك له بين التراكيب الهيدروجينية؟
  \item كم \omterm{def:b2:diatomic-mos:lcao}{مدارًا جزيئيًا} ينتج، وكم منها ممتلئ؟
\end{enumerate}

\medskip
\noindent\textbf{الجزء الثاني --- H--O--H خطي.}
\begin{enumerate}[resume]
  \item في جزيء خطي على طول $z$، أي تركيب يتآثر مع
        المدار $2\mathrm p_z$ للأكسجين؟
  \item وأيهما يتآثر مع المدار 2s للأكسجين؟
  \item أي مدارات الأكسجين تبقى غير رابطة، وبأي تعدد للتساوي في الطاقة؟
  \item املأ المستويات بإلكترونات التكافؤ الثمانية.
  \item هل كان الماء الخطي سيكون بارامغناطيسيًا؟
  \item أين كان سيقع أعلى مستوى مشغول له، مقارنة بمدار 2p
        للأكسجين؟
\end{enumerate}

\medskip
\noindent\textbf{الجزء الثالث --- الانحناء.}
\begin{enumerate}[resume]
  \item يُحنى الجزيء الخطي الممتد على طول $z$ في المستوي $yz$. أي
        \omterm{def:b2:diatomic-mos:bonding}{مدار غير رابط} يبدأ بالتداخل مع $h_1 + h_2$؟
  \item ماذا يحدث للمستوى الممتلئ الذي يحمله؟
  \item أي مستوى رابط يرتفع قليلًا، ولماذا؟
  \item استنتج: لماذا يكون الماء منحنيًا؟
  \item قارن الزاوية المقيسة بالزاوية $90^\circ$ (ارتباط p خالص) و$109.5^\circ$.
  \item أي \omterm{def:b2:diatomic-mos:bonding}{مدار غير رابط} لا يمسّه الانحناء؟
\end{enumerate}

\medskip
\noindent\textbf{الجزء الرابع --- طيف الإلكترونات الضوئية.}
\begin{enumerate}[resume]
  \item كم حزمة تكافؤ يُظهر الطيف؟
  \item بتقريب كوبمانز، من أي مدار تأتي الحزمة الأولى؟
  \item لماذا تكون هذه الحزمة ضيقة، بينما تكون حزم \omterm{def:b2:diatomic-mos:bonding}{المدارات الرابطة} عريضة؟
  \item قارن طاقة التأين الأولى للماء بطاقة تأين ذرة
        الأكسجين، وفسّر الفرق.
  \item ماذا يحل «بالزوجين الحرين المتكافئين» في بنية لويس؟
  \item اذكر عدد حزم التأين التكافؤية للماء وطاقة
        أدناها طاقةً.
\end{enumerate}
\end{problem}
